\section{Decremento y conquista}

El paradigma de \concept{decremento y conquista} aprovecha la relación entre el problema y una versión más pequeña del problema, que le llamamos \concept{subproblema}. La idea es imaginar que ya se conoce la solución a ese subproblema y aprovechar eso para terminar de resolver el problema.

El subproblema se obtiene reduciendo el problema a distintas tasas:
\begin{itemize}
    \item \concept{Decremento por una constante}: se reduce el tamaño del problema en una cantidad fija en cada iteración del algoritmo, usualmente en uno.
    \item \concept{Decremento por un factor constante}: se reduce el problema en una proporción fija en cada iteración, usualmente a la mitad.
    \item \concept{Decremento de tamaño variable}: se reduce el tamaño en distinta cantidad en cada iteración.
\end{itemize}

Se presentan dos algoritmos para concretar esta idea, correspondientes a los dos primeros tipos. Deliberadamente de omite un ejemplo que reduzca el problema en valores variables; uno clásico es el Algoritmo de Euclides, que calcula el máximo común divisor de dos enteros, pero su estudio requiere de varias nociones de teoría de números.

\setinputbase{diseno/decremento_conquista}
\inputlist{insertion_sort, busqueda_binaria, refs}
